%\section{The frequent random lookup uses distributed nonlinear computation}
\section{Distributed nonlinear computation}
\label{sec:mechanism}

With trained weights fixed, replacing GELU with the identity reduces
random-relation accuracy from 100\% to 51\%. Making attention uniform or
removing all attention/MLP outputs also reduces accuracy to near chance
(Fig.~\ref{fig:goal2-mechanism}B). These interventions show that the trained
lookup depends on nonlinear activation, attention routing, and MLP computation.


The random answer becomes reliably decodable only in the final block
(Figure~\ref{fig:goal2-mechanism}A). Yet removing any single attention or MLP
sublayer reduces accuracy, indicating that earlier blocks also contribute to the
lookup (Figure~\ref{fig:goal2-mechanism}D).

Residual-stream patching shows a progressive shift in where interventions recover
the clean answer margin (Figure~\ref{fig:goal2-mechanism}C). Recovery decreases
across blocks when patching the subject position and increases when patching the
object position. This pattern is consistent with subject information being
transferred to the object position, where the model produces its answer.

These interventions impair both relation types, so they do not explain the
difference in affine decodability. That distinction comes from the raw-embedding
OLS comparison.

